\section{Independent verification of the full finite prime-orbit
relation
theorem}\label{independent-verification-of-the-full-finite-prime-orbit-relation-theorem}

24 September 2026. Mathematical check FPC0--FPC6.

\subsection{FPC0. Reading coverage and unchanged
source}\label{fpc0.-reading-coverage-and-unchanged-source}

The complete \href{FAITHFUL_GLOBAL_PRIME_ORBIT.md}{FPO0--FPO8} was read
and checked against
\href{GLOBAL_GAUSSIAN_ORIGINAL_ZETA_TRACE.md}{GGT0--GGT8}, the original
PGD multiplier comparison, and the actual ASD14 cone. Source
\(Z_0,Z_1/\tau,Z_2\), all retractions, and the receiving-operation
domains are unchanged. In particular the arithmetic constant function
\(1\) is not primitive \(\tau\).

The original author TeX of Alain Connes and Caterina Consani,
\href{https://arxiv.org/abs/2006.13771v1}{Weil positivity and Trace
formula, the archimedean place, arXiv:2006.13771v1}, was read at
Appendix ``Explicit formula,'' lines2033--2069, and the immediately
following Proposition mainprop and its proof. The bibliography entry for
Enrico Bombieri's The Riemann hypothesis (2006), credited by that
appendix, was checked. No complete reading of Bombieri's chapter is
claimed.

The exact author source was
weil-compo.tex
SHA256 \nolinkurl{b01d353b0423b6fedee373c3c33fe3678eea733f62049810750f6c64ef20f3fc}.
The appendix's ``all complex zeros'' wording at line2042 is read with
its explicit nontrivial-zero convention in Proposition mainprop,
line2075, and the original contour calculation in GGT. An extra infinite
series over negative even integers is not inserted into that formula.

\subsection{FPC1. The compact-support extension retains every
term}\label{fpc1.-the-compact-support-extension-retains-every-term}

For the actual Gaussian \[
f(u)=\frac{T}{2\sqrt\pi}
e^{-T^2\log^2(u/x)/4},\qquad h(v)=f(e^v),
\] take the stated smooth cutoff \(h_R(v)=\eta(v/R)h(v)\), with
\(\eta=1\) on \([-1,1]\) and zero outside \([-2,2]\). Every derivative
of \(h\) is a fixed polynomial in \(v-\log x\) times its Gaussian. The
difference \(h_R-h\), with every derivative, is supported in
\(|v|\ge R\). The product rule therefore proves \[
\|e^{A|v|}(h_R-h)^{(k)}\|_{L^1}
+\|e^{A|v|}(h_R-h)^{(k)}\|_\infty\longrightarrow0
\] for every fixed \(A,k\). Derivatives falling on the cutoff contribute
only bounded factors \(R^{-j}\), and a Gaussian tail dominates all the
resulting fixed exponential-polynomial weights.

Integration by parts in the Mellin transform on logarithmic coordinates
proves convergence in every \(\mathcal B\) strip seminorm. The zero
count gives \(\sum_\rho m_\rho(1+|\Im\rho|)^{-3}<\infty\), so this
implies convergence of the full nontrivial-zero trace. Endpoint
evaluations at zero and one are continuous on \(\mathcal B\).

The two prime evaluations obey \[
|f_R(n)-f(n)|\le\varepsilon_R n^{-A},\qquad
|f_R^\sharp(n)-f^\sharp(n)|\le\varepsilon_R n^{-A-1},
\quad \varepsilon_R\to0.
\] For \(A>2\), the sums weighted by \(\Lambda(n)\le\log n\) converge.
Thus every prime repetition passes to the limit.

For the original archimedean expression, logarithmic substitution gives
exactly \[
\mathcal W_{\mathbb R}(f)
=(\log4\pi+\gamma)h(0)
+\int_0^\infty
\frac{h(v)+e^{-v}h(-v)-2e^{-v}h(0)}
{1-e^{-2v}}\,dv.
\tag{FPC1.1}
\] Its numerator vanishes at zero and has a bounded derivative there. In
the difference for \(h_R-h\), the complete numerator is identically zero
near zero and its value term is exactly zero. On the complement of a
fixed neighborhood of zero the denominator is bounded below, and the
weighted \(L^1\) convergence proves convergence of the whole integral.
This verifies the cutoff extension used in FPO3, including its singular
endpoint and reflected term.

\subsection{FPC2. Direct independent Gamma-integral
comparison}\label{fpc2.-direct-independent-gamma-integral-comparison}

There is a direct verification of the archimedean equality that does not
require subtracting the two explicit-formula identities. Write \[
\Phi(s)=x^se^{s^2/T^2},\qquad
g_\chi(s)=\frac{\chi'(s)}{\chi(s)},\qquad
\chi(s)=\pi^{s-1/2}\frac{\Gamma((1-s)/2)}{\Gamma(s/2)}.
\] For \(\Re z>0\), the Euler digamma series used in GGT4 gives \[
\psi(z)+\gamma
=\int_0^\infty\frac{e^{-u}-e^{-zu}}{1-e^{-u}}\,du.
\tag{FPC2.1}
\] Indeed expand the denominator geometrically and integrate the paired
terms. Their absolute integral sum converges: on each pair the
mean-value integral bounds the difference by a constant times
\(u e^{-(n+c)u}\), with \(c=\min(1,\Re z)>0\), whose integral is
\(O((n+c)^{-2})\). This justifies interchange and recovers the
convergent digamma difference series.

On the actual middle edge \(s=1/2+it\), the two digamma arguments have
real part \(1/4\). Equation (FPC2.1) gives \[
g_\chi(1/2+it)=\log\pi+\gamma
+\int_0^\infty
\frac{e^{-u/4}\cos(tu/2)-e^{-u}}{1-e^{-u}}\,du.
\tag{FPC2.2}
\] This retains the full original Gamma quotient.

The double integral against \(\Phi(1/2+it)\,dt/(2\pi)\) is absolutely
integrable. For \(0<u\le1\), the absolute value of the quotient in
(FPC2.2) is bounded by \(C(1+t^2u)\), using
\(|\cos(tu/2)-1|\le t^2u^2/8\). For \(u\ge1\) it is bounded by
\(Ce^{-u/4}\). The Gaussian in \(t\) makes both bounds integrable. Hence
Fubini applies before Fourier inversion.

Original Mellin inversion, with its unshifted definition, gives \[
\frac1{2\pi}\int_{\mathbb R}\Phi(1/2+it)\,dt=h(0),
\] \[
\frac1{2\pi}\int_{\mathbb R}
\Phi(1/2+it)\cos(tu/2)\,dt
=\frac12\left[e^{-u/4}h(-u/2)+e^{u/4}h(u/2)\right].
\tag{FPC2.3}
\] Substitution into (FPC2.2), followed by the exact variable change
\(u=2v\), proves \[
A_T^{(1/2)}(x)
=(\log\pi+\gamma)h(0)
+\int_0^\infty
\frac{h(v)+e^{-v}h(-v)-2e^{-2v}h(0)}
{1-e^{-2v}}\,dv.
\tag{FPC2.4}
\] All integrals have their ordinary convergent endpoint values. Compare
this with the author's exact expression (FPC1.1). Their difference is \[
\begin{aligned}
A_T^{(1/2)}(x)-\mathcal W_{\mathbb R}(f)
&=-\log4\,h(0)
+2h(0)\int_0^\infty
\frac{e^{-v}-e^{-2v}}{1-e^{-2v}}\,dv\\
&=-\log4\,h(0)+2\log2\,h(0)=0.
\end{aligned}
\tag{FPC2.5}
\] The middle integral is exactly
\(\int_0^\infty e^{-v}/(1+e^{-v})\,dv=\log2\). This checks the
coefficient \(\log4\pi+\gamma\), including the tail subtraction in the
source distribution.

Finally the full logarithmic derivative \(g_\chi\) has residue \(+1\) at
zero. The finite shift from \(\Re s=-1\) to \(1/2\) therefore gives \[
A_T^{(1/2)}(x)-A_T(x)=\Phi(0)=1.
\] Together with (FPC2.5), this proves independently \[
\boxed{\mathcal W_{\mathbb R}(f_{T,x})=A_T(x)+1.}
\tag{FPC2.6}
\] The value \(1\) is the original Mellin endpoint integral, not the
Gaussian's point value \(f(1)\), and not primitive \(\tau\).

\subsection{FPC3. The quotient really is an algebra, while the full
cokernel is a
module}\label{fpc3.-the-quotient-really-is-an-algebra-while-the-full-cokernel-is-a-module}

The multiplier algebra \(\mathcal M\) in FPO1 is closed under
multiplication. If \(H\in\mathcal M\), \(F\in\mathcal B\), then \[
b_{A,N}(HF)\le C_A b_{A,N+d_A}(F),
\] so \(\mathcal B\) is an ideal of \(\mathcal M\). The full-zero-jet
condition defining \(\mathcal J\) is preserved by multiplication, by the
finite Leibniz formula at every actual zero. Thus
\(\mathcal B+\mathcal J\) is an ideal, and
\(\mathcal A_\zeta=\mathcal M/(\mathcal B+\mathcal J)\) is a
well-defined quotient algebra.

The map from the finite group algebra sends \[
\delta_r\mapsto[r^s].
\] The exact equality \(r^st^s=(rt)^s\), using the real logarithms of
positive numbers, proves multiplicativity. Finite sums preserve this
identity by expansion, and the arithmetic unit goes to the arithmetic
unit. This verifies an algebra map into \(\mathcal A_\zeta\).

PGD's map \(\mathcal A_\zeta\to\mathcal C_\zeta\) is proved linear and
injective. No multiplication on all of \(\mathcal C_\zeta\) is part of
that construction. FPO5.8 correctly keeps this distinction. The
additional intertwining \[
T_a^\diamond[\lambda_{r^s}]=[\lambda_{(ar)^s}]
\] identifies the injected space with the regular group module. It does
not require an algebra structure on the complete cokernel.

\subsection{FPC4. Full zero-jet relations imply the decisive value
relations}\label{fpc4.-full-zero-jet-relations-imply-the-decisive-value-relations}

A zero class in \(\mathcal A_\zeta\) means that for some
\(F\in\mathcal B\), the difference \[
H(s)-F(s),\qquad H(s)=\sum_{i=1}^k a_i r_i^s,
\] has all actual required zero jets equal to zero. In particular
\(H(\rho)=F(\rho)\) at every actual zero. Using only this consequence is
valid: the proof is not weakening the ideal used to define the quotient.

For a fixed support index \(j\), all sums in \[
\sum_{i=1}^k a_i S_T(r_i/r_j)
=\sum_\rho m_\rho r_j^{-\rho}F(\rho)e^{\rho^2/T^2}
\] converge absolutely. The right side is bounded uniformly for
\(T\ge1\) by \[
e\max(1,r_j^{-1})\,b_{1,3}(F)
\sum_\rho m_\rho(1+|\Im\rho|)^{-3}.
\tag{FPC4.1}
\] This finite bound uses the original quotient representative and
retains the actual zero multiplicities.

Exactly one left-side scale is \(1\), since the listed \(r_i\) are
distinct. GGT gives \[
\sum_{i=1}^k a_i S_T(r_i/r_j)
=\frac{a_j}{2\sqrt\pi}T\log T+O(T).
\tag{FPC4.2}
\] The finite sum of all other fixed-scale errors is \(O(T)\), including
every prime-power resonance. Dividing the equality by \(T\log T\) gives
\(a_j=0\). This works for complex coefficients as well: the absolute
value of the remainder divided by \(T\log T\) tends to zero, and the
nonzero real scalar \(1/(2\sqrt\pi)\) multiplies \(a_j\). Repeating for
each \(j\) proves that the finite relation is zero.

This verifies the injectivity in the actual quotient, which is stronger
than linear independence of the original exponential functions before
quotienting.

\subsection{FPC5. Prime algebra, actual cone cycles, and support
labels}\label{fpc5.-prime-algebra-actual-cone-cycles-and-support-labels}

Unique integer factorization proves that \[
\bigoplus_{p\ {\rm prime}}\mathbb Z
\longrightarrow\mathbb Q_+^\times,\qquad
(n_p)\longmapsto\prod_p p^{n_p}
\] is an isomorphism of groups. Only finite-support exponent tuples
occur. Thus the corresponding finite group algebra is exactly the
Laurent polynomial algebra in all prime variables, with each polynomial
involving finitely many variables and monomials. This verifies FPO6.1
and its use of all negative and positive prime exponents.

The actual ASD14 cycle associated with a multiplier is
\(\pi'\lambda_H\). Since \(\pi d=0\), its image under \(d'\) is zero. If
this cycle is a boundary, then \(\pi'\lambda_H=\pi'D\pi(b)\).
Injectivity of \(\pi'\) gives \(\lambda_H\in DQ\), the exact relation
excluded by FPC4 for a nonzero finite group-algebra element. Thus the
orbit injection survives in the actual cone cohomology, with the stated
positive middle differential.

The lift \((v,\alpha)\mapsto(f(v),\alpha)\) is well-defined on the
existing support carrier because nonzero amplitudes occur only at top
support and a zero output retains its label. Injectivity of the original
receiving map implies injectivity of its lift. The separate faithful
closed source copies remain in the full complex and are not removed by
the orbit calculation.

\subsection{FPC5A. Verification of the appended supported-cone
propagation}\label{fpc5a.-verification-of-the-appended-supported-cone-propagation}

The final FPO8 paragraph now invokes UOS3--UOS8A in
\href{GLOBAL_UNIT_ORBIT_AND_SUPPORTED_COMPARISON.md}{the complete
supported comparison}. I read its original localization row, transpose
row, chain comparison, cone and cohomology maps in UOS3--UOS6, and the
source actions, Fourier section and full orbit propagation in
UOS8--UOS8A.

Write \(Y=\chi\otimes Q'\), \(M_\zeta=D_\zeta Q\), and
\(\alpha\lambda=(\lambda,-\lambda)\). The actual degree-one map of cones
sends the cycle \(\pi'\lambda\) to \((\pi'\lambda,-\pi'\lambda)\). The
displayed target cohomology is \(Y^2/\alpha M_\zeta\), and its exact
decomposition is \[
[(\lambda_+,\lambda_-)]\longmapsto
\left(\left[\frac{\lambda_+-\lambda_-}{2}\right],
\frac{\lambda_++\lambda_-}{2}\right)
\in (Y/M_\zeta)\oplus Y.
\] Substituting \((\lambda,-\lambda)\) gives \(([\lambda],0)\). A change
\(\lambda\mapsto\lambda+D_\zeta q\) changes the representative by
\(\alpha D_\zeta q\), which is precisely the subspace divided out. Thus
the appended map is exactly \(c_r\mapsto(c_r,0)\), and is injective
without assuming a splitting \(Q'\to Q\).

In degree \(-1\), the chain map is the identity on \(P\); its
restriction to \(\ker d=H\) is the identity. In degree \(2\), the chain
map is the identity on \(\chi P'\); the resulting map from \(\chi H'\)
to \(\chi E'\) is restriction along \(E\hookrightarrow H\). In the
original coordinates it keeps the four endpoint coordinates and both
extra closed dual copies and removes only the \(S'\) coordinate. This
verifies the map stated in FPO8, including its direction.

Finally UOS8's original Fourier section has restriction equal to
identity by \(\Sigma\widehat h=R\Sigma h\), with the retained minus
extra coordinate and zero plus extra coordinate. Its localization
connecting map is consequently zero. The supported placement of the
residue-comparison classes does not replace that section or contradict
its existence. No additional hypothesis is needed for the appended
propagation.

\subsection{FPC6. Outcome and the identified typesetting
defect}\label{fpc6.-outcome-and-the-identified-typesetting-defect}

FPO0--FPO8's mathematical claims checked above are valid: the
compact-support extension converges in every required term; the full
archimedean comparison is exactly \(A_T+1\); the group-algebra map is
injective after the actual multiplier quotient; the map into the
complete cokernel is linear and equivariant; and the full-zero-jet
relation argument retains every multiplicity.

One typesetting defect was found in the reviewed Markdown: FPO4.3,
line166, contained two literal tab characters followed by ``frac12,'' in
place of the two commands \(\backslash\mathrm{tfrac}12\). Both
backslashes in the two one-half digamma coefficients have now been
restored by the parent author. This does not change the verified
mathematical signs or factors. The check did not overwrite their file.

The result concerns the exact original receiving comparison and its
finite orbit relations. It does not identify the complete cokernel with
an algebra, assert a purity bound, or identify a nonzero cokernel class
with an RH counterexample.
